%%==================================================
%% chapter01.tex for Doctoral Thesis
%% Chapter 1: Introduction
%% Narrative: training goal -> two core gaps -> environment foundation + experimental verification
%%==================================================

% !TeX root = ../main.tex

\chapter{绪论}
\enchapter{Introduction}
\label{chap:introduction}

\section{研究背景和意义}
\ensection{Background and Significance}

化学实验是化学教学、科学研究以及工业生产中的重要基础环节。无论是基础化学规律的验证，还是新材料、新药物及新工艺的探索，都离不开规范、可靠的实验操作过程。然而，传统化学实验通常以人工操作为主，实验过程对人员经验、操作熟练度和规范执行能力依赖较强，存在较高的人工成本。此外，化学实验中常涉及腐蚀性、易燃性或有毒试剂，一旦操作不当，容易引发安全风险，造成严重的后果。因此，借助自动化、数字化与智能化技术降低人工负担、降低实验风险并提升实验效率，已成为化学实验相关研究的重要发展方向。

随着人工智能、机器人技术与虚拟仿真技术的快速发展，具身智能为化学实验自动化提供了新的解决路径。具身智能强调智能体在具体环境中的感知、决策与执行闭环，适合处理化学实验这类同时涉及空间理解、对象操作、流程执行与结果反馈的复杂任务。对于化学实验这一安全敏感且操作规范性强的任务场景而言，将具身智能与化学实验相结合，具有明显的研究价值和应用前景。

然而，面向化学实验的具身智能技术有着独特的挑战。首先，化学实验场景具有较强的专业性与安全约束，其空间布局、物体摆放和操作流程不仅要满足一般物理可行性，还需要符合实验规范与安全要求，这与日常室内场景存在明显差异。其次，现有开源三维资产库主要面向日常家具、生活用品等通用对象，化学实验所需的大量专业物体资产严重匮乏，尤其是烧杯、量筒等透明物体和小物体，难以直接从现有开源数据集中获得高质量三维资产。最后，现有研究对于化学实验场景下从实验手册文本到机器人可执行操作序列的转化关注不足。

近年来，随着大语言模型与三维重建技术的发展，为解决上述问题提供了新的技术机遇。一方面，大语言模型在自然语言理解、知识推理和结构化生成方面展现出较强能力，使其能够更好地理解实验文本中的操作语义、对象关系与流程逻辑，为复杂实验场景构建和任务表示生成提供支持。另一方面，以  NeRF\cite{mildenhall2020nerf}  为代表的神经隐式场算法在三维几何建模与外观重建方面展现出较好的表现，为高保真化学资产重建提供了新的实现路径。这些技术的发展，使得面向化学实验的具身智能研究逐步走向可实现、可验证的阶段。

基于上述背景，围绕化学实验场景开展具身智能关键技术研究，具有明确的理论意义与应用价值。从理论层面看，该方向有助于推动大语言模型、三维场景生成与具身智能执行等技术在化学实验场景中的交叉融合，为自然科学专业场景下的具身智能研究提供新的问题牵引与方法范式。从应用层面看，相关研究可为智能实验室、虚拟实验教学和数字孪生实验系统提供基础支撑，并为后续更复杂的实验机器人训练与验证奠定基础。综上所述，本文围绕面向化学实验的具身智能关键技术展开研究，重点关注三维场景生成、高保真化学资产重建以及具身智能任务规划这三个核心问题。



\section{国内外研究现状}
\ensection{Related Work}


\subsection{三维场景生成算法}

三维场景生成是计算机图形学与人工智能领域的重要研究问题，其目标是在给定空间范围内自动生成具有合理结构、语义关系和物体布局的三维环境。早期的研究~\cite{germer2009procedural,merrell2011furniturelayout,yu2011makeithome,wonka2003instant}主要依赖人工规则、模板组合与程序化生成方式，通过预定义的布局规则、语法结构生成室内场景。这类算法通常具有较强的可控性，能够在一定程度上保证场景结构的基本合理性，但其生成结果较大程度依赖于人工设计的规则库，难以适应复杂多变的空间关系，也较难覆盖真实场景中的丰富布局模式。

受益于大规模三维数据集的规模化增长\cite{song2017sscnet,song2015sunrgbd,fu2021future,fu2021front}，研究者开始采用统计学习算法对室内场景分布进行建模。相关工作~\cite{wang2019planit,li2019grains,purkait2020sgvae,gao2023scenehgn}通过分析真实场景中物体的共现关系、相对位置、朝向与功能分区等统计规律，实现更加数据驱动的场景合成。这类算法相比纯规则算法减少了人工设计负担，使生成场景在整体布局上更接近真实数据分布，但在复杂语义约束、多目标优化以及跨场景泛化能力方面仍然存在一定局限。

近年来，基于深度学习的场景生成算法逐渐成为主流\cite{wang2021sceneformer,zhang2022fast3d,wei2023legonet,zhai2023commonscenes,yang2024physcene,paschalidou2021atiss,tang2024diffuscene,lin2024instructscene}。代表性算法如 Atiss\cite{paschalidou2021atiss} 和 Diffuscene\cite{tang2024diffuscene}，通过自回归模式或扩散生成机制学习室内场景的布局分布，能够自动生成在视觉上更自然、在统计上更符合真实数据的三维房间布局。在此基础上，Lin 等人的工作\cite{lin2024instructscene}进一步引入文本条件、语义关系图或用户指令，以增强生成结果的可控性与语义一致性。这些算法显著提升了三维场景生成的自动化程度，使得场景生成从“规则驱动”逐步转向“数据驱动”和“语义驱动”。

随着大语言模型的发展，研究者进一步尝试将自然语言理解能力引入场景生成任务，通过大语言模型直接生成布局方案、空间关系或功能分区，并结合约束求解器或检索机制提升生成结果的可执行性与一致性\cite{aguinakang2024openuniverse,celen2024idesign,hu2024scenecraft,ocal2024sceneteller,yang2024llplace,wang2024chat2layout,feng2023layoutgpt,yang2024holodeck}。这些工作表明，大语言模型在理解高层用户意图、生成语义布局结构以及辅助场景自动构建方面展现出良好潜力。这类算法为具身智能场景生成提供了新的思路，也使场景生成从“学习数据分布”进一步走向“理解任务需求与用户意图”。

对于化学实验场景生成还需要考虑化学品安全操作规程与空间布局约束之间的对应关系。Shimizu 等人的工作\cite{shimizu2018pattern} 面向化学实验过程提出了基于模式的本体表示方法，用于描述化学性质、物质状态、实验条件、操作动作以及危害分类等信息。Zheng 等人的工作\cite{zheng2021hazardouskg} 进一步面向危险化学品管理构建知识图谱，将危险化学品属性、事故记录和管理规则组织为可查询、可推理的结构化表示。这些工作表明，化学品安全操作规程可以被转化为结构化知识并支持风险识别，但在三维场景生成中仍需进一步将安全属性和操作规程，显式建模为约束，以保证生成场景的安全性。

尽管现有三维场景生成算法已经能够较好地建模室内场景的统计规律与语义关系，但其主要应用对象仍然是家庭、办公等通用场景，优化目标也大多集中在视觉合理性、统计一致性和语义匹配上。对于面向具身智能的化学实验场景而言，更重要的是实验的安全性与机器人的可操作性。例如，危险化学品与热源之间的隔离要求、实验室是否能满足操作需求，均属于必须实现的目标。因此，如何生成一个安全、可操作、高质量的化学实验场景，仍然是一个重要问题。


\subsection{三维物体重建算法}
\label{subsec:asset_reconstruction}

传统三维资产重建方法主要建立在多视图几何理论基础之上，其核心思想是利用不同视角下的图像观测恢复物体的三维结构。围绕这一技术路线，研究者提出了一系列经典方法。例如，MVS\cite{furukawa2010mvs} 通过结合局部光度一致性与全局可见性约束，对初始稀疏点进行逐步扩展，从而获得更加稠密的三维点云。Poisson \cite{kazhdan2006poisson} 则为稠密点云向连续三角网格的转换提供了经典方案。此后，COLMAP\cite{schonberger2016sfm} 进一步整合了传统几何重建中的关键步骤，逐步成为该类方法的代表性范式。具体而言，COLMAP 通常首先通过 Structure-from-Motion 从无序图像中恢复相机外参与稀疏点云，建立场景的初始几何结构，随后结合 Multi-view Stereo 对图像进行逐像素深度估计，将稀疏结构进一步稠密化并通过表面重建方法生成显式三角网格，并可进一步结合纹理映射得到较完整的三维资产。总体来看，这类方法在理论和工程实现上都较为成熟，COLMAP 所体现的"位姿恢复—稠密重建—表面生成"流程也逐渐成为传统几何重建中的经典技术框架，并在不透明、纹理较丰富且成像条件相对稳定的物体上取得了较好的重建效果。


近年来，基于神经辐射场的三维重建算法为复杂物体建模提供了新的研究思路。NeRF\cite{mildenhall2020nerf} 首次将连续空间中的体积密度与视角相关辐射度统一建模，并结合可微体渲染实现了高质量的新视角合成与三维表示学习，为后续神经隐式重建方法奠定了基础。围绕原始 NeRF 在场景范围、渲染质量、计算效率以及输入视图依赖性等方面的局限，研究者提出了一系列改进方法。NeRF++\cite{zhang2020nerfplusplus} 通过改进场景参数化方式提升了模型对无界场景的表示能力，mip-NeRF\cite{barron2021mipnerf} 及 mip-NeRF 360\cite{barron2022mipnerf360} 通过多尺度表示与抗锯齿建模增强了不同尺度和大范围场景下的渲染稳定性，Ref-NeRF\cite{verbin2022refnerf} 进一步提升了对反射表面的视角相关外观建模能力，Instant-NGP\cite{muller2022instantngp} 则利用多分辨率哈希编码显著降低了训练与渲染开销。另一方面，为了提升稀疏视图条件下的泛化能力，研究者还从跨场景先验学习角度展开探索，PixelNeRF\cite{yu2021pixelnerf}、IBRNet\cite{wang2021ibrnet} 和 MVSNeRF\cite{chen2021mvsnerf} 等方法，通过引入图像特征投影、多视图特征融合和几何代价体表示，增强了弱观测条件下的三维表达能力。尽管 NeRF 类方法在视图合成和连续场表示方面展现出显著优势，但其几何表面通常依赖体密度场的间接提取，难以稳定恢复高精度显式表面。为此，研究者进一步将辐射场建模与隐式表面表示相结合，提出了 UNISURF\cite{oechsle2021unisurf}、VolSDF\cite{yariv2021volsdf} 和 NeuS\cite{wang2021neus} 等方法。其中，NeuS 通过将有符号距离函数与体渲染过程紧密耦合，显著提升了多视图条件下的表面几何恢复质量。

尽管传统多视图几何方法与 NeRF 系方法均推动了三维重建技术的发展，但在面向化学实验的专业三维资产重建时仍存在明显不足。对于传统方法而言，其基于显式几何和多视图匹配的流程在纹理映射与显式网格生成方面较为成熟，因此在不透明、纹理丰富的物体上通常能够获得较好的重建结果。然而，化学实验中常见的烧杯、试管、量筒和冷凝管等物体往往具有透明玻璃材质与小尺度结构，容易导致深度估计不稳定、跨视图匹配失败以及局部细节缺失。相比之下，NeRF 类方法通过连续隐式表示与体渲染机制，在透明材质建模方面展现出更大潜力，为化学实验物体重建提供了新的技术路径。但对于面向仿真平台和具身智能系统使用的三维资产而言，仅有隐式表示仍然是不够的。因此，如何面向透明和小尺度化学实验物体实现高质量重建，并进一步完成从隐式表示到显式三维资产的稳定转化，是该方向需要解决的关键问题。


\subsection{面向化学实验的具身智能任务规划}
\label{subsec:embodied_planning}

在具身智能任务中，机器人根据自然语言命令生成可执行的动作序列是核心任务之一。随着大语言模型在语义理解与逻辑推理方面展现出卓越能力，研究者开始探索将其作为高层任务规划器与指令生成模块。SayCan\cite{ichter2023saycan} 通过结合大语言模型输出与动作可执行性评分实现长程任务规划，Huang 等人的工作\cite{huang2023innermonologue} 进一步引入环境反馈的闭环推理，Liang 等人的工作\cite{liang2023codeaspolicies} 则将自然语言任务转化为可调用感知结果和控制接口的策略代码。PaLM-E\cite{driess2023palme} 将视觉感知、语言理解与机器人控制进一步统一到同一模型框架中，增强了机器人在真实环境中的多模态任务理解能力，RT-2\cite{zitkovich2023rt2} 则进一步探索了由视觉语言知识向机器人动作控制迁移的可能性。与此同时，VoxPoser\cite{huang2023voxposer} 通过结合大语言模型与可组合的三维价值图，将自然语言指令转化为可执行的操作规划。这些工作表明，大语言模型在具身智能系统中适合作为高层语义规划模块，负责理解任务目标、组织执行流程并调用已有技能。

与此同时，大语言模型在预训练阶段已经从大规模语料中积累了大量化学相关知识，展现出一定的化学专业理解与推理能力，因此相关研究开始探索将大语言模型引入化学问题求解、实验规划与流程组织之中。ChemCrow\cite{bran2024chemcrow} 通过将大语言模型与化学工具库结合，增强了模型在合成规划、分子设计和材料任务中的求解能力。Coscientist\cite{boiko2023coscientist} 进一步展示了大语言模型结合工具调用、文档检索和实验自动化平台后在实验规划与执行中的潜力。这些工作表明，大语言模型不仅能够用于化学知识推理，也能够在实验目标理解、流程组织与实验自动化系统交互中发挥重要作用。


尽管现有研究已经分别在通用具身智能任务规划和大语言模型辅助化学任务求解方面取得了良好进展，但面向化学实验的具身智能任务规划研究仍然相对缺乏。一方面，现有具身智能规划方法大多面向家庭服务、导航或一般操作任务，对化学实验中安全规范、实验步骤考虑不足，难以直接适用于化学实验场景。另一方面，现有化学领域相关工作虽然证明了大语言模型在化学知识推理和实验流程组织中的潜力，但其研究重点更多集中于化学问题求解、工具调用与实验方案生成，对于面向化学实验的具身智能任务规划仍研究不足。



\section{本文主要内容}
\ensection{Main Content}
\label{sec:intro_main_contributions}

本文针对上述问题，围绕基于大语言模型的化学实验具身智能技术，从化学实验场景生成、高保真化学资产重建和面向化学实验的具身智能任务规划三个方面开展研究。

（1）针对场景生成算法在化学安全约束和多样性等方面的不足，本文提出了由大语言模型驱动的化学实验场景生成算法，生成安全、高质量和多样化的化学实验场景。具体而言，本文首先设计了面向大语言模型的三维资产检索方案，增强大语言模型对三维资产的语义理解能力。随后，利用大语言模型对自然语言实验描述进行解析，生成物体清单与实验步骤，并对物体清单进行安全补齐。在此基础上，将实验流程中的交互关系、安全规范和空间布局要求统一表示为物体约束图，通过求解约束图得到实验场景。考虑到复杂实验场景中人工自由浏览成本大、效率低，本文进一步研究了面向场景展示的相机规划算法，通过综合考虑关键区域显著性、场景覆盖率和相机运动平滑性，生成高效展示场景显著区域的相机路径。最后，为了进一步提升场景质量并增强训练环境的多样性，本文提出面向具身智能训练的场景进化算法，在安全约束基础上，结合大语言模型驱动的交叉与变异算子以及人工反馈，迭代生成兼顾质量与差异性的场景集合。实验结果表明，本文方法能够生成安全、高质量且多样化的化学实验场景，为具身智能体训练提供有效的环境支撑。

（2）针对化学实验资产类型多样问题，本文先提出了化学实验物体分类体系图，针对各类别物体重建所带来的挑战，本文研究了基于隐式神经场的高保真化学资产重建算法。具体而言，本文以 NeuS 为基础重建框架，针对小物体在输入图像中像素占比较低、训练早期监督信号稀疏的问题，本文提出基于学习态增强的小物体重建加速策略，通过动态调节采样方向与距离场梯度方向的作用强度，加速训练前期的几何收敛过程，使网络更早进入细节学习阶段，从而提高了几何模型细节重建质量。针对透明物体表面与零等值面不稳定对应的问题，采用面向透明物体的统一几何表面提取策略，通过在绝对距离场上搜索局部极小值，实现透明区域与不透明区域表面的统一恢复，从而提升透明物体的几何提取质量。在此基础上，为满足主流仿真平台对显式模型的使用需求，本文研究了基于 NeuS 的高分辨率纹理贴图生成算法，利用 UV 参数化、光栅化烘焙、形态学修复与 KNN 补全等步骤，将隐式颜色场转换为高质量显式纹理资产。实验结果表明，本文方法能够生成高精度的几何模型与高分辨率的纹理贴图，从而为仿真平台、机器人感知提供高保真的三维资产基础。

（3）针对面向化学实验的具身智能任务规划问题，本文研究了从实验手册文本到操作原语的分层转化方法，以支持机器人在仿真环境中执行化学实验流程。具体而言，本文采用“实验手册文本—步骤计划—操作原语”的分层表示思路，首先利用大语言模型对实验手册进行解析，提取实验过程中的关键操作、涉及对象及其参数要求，并生成步骤计划。随后，结合预定义的机器人操作原语集合，将步骤计划进一步编译为机器人控制器可调用的操作原语序列。最后，本文在 NVIDIA Omniverse 机器人仿真平台中构建化学实验执行环境，并通过经典无机化学反应案例展示和操作原语生成质量与成本的定量实验，对所提方法进行了系统验证。实验结果表明，本文方法能够较为准确地完成从实验手册文本到机器人操作原语序列的转化，并在可执行性与总体生成成本方面表现出较好的综合性能，为面向化学实验自动化的具身智能机器人系统提供了可行的任务规划框架。


\section{章节安排}
\ensection{Organization}
\label{sec:intro_organization}

\begin{figure}[htbp]
    \centering
    \includegraphics[width=0.9\textwidth]{figures/ch1/ch_framework.png}
    \caption{本文技术路线与章节组织结构图。}
    \label{fig:ch_framework}
\end{figure}

图~\ref{fig:ch_framework} 展示了本文的整体技术路线与章节组织结构。该图从化学实验场景生成、高保真化学资产重建以及机器人任务规划三个层面，概括了本文围绕虚拟化学实验环境构建所开展的主要研究内容。基于这一技术路线，本文各章节安排如下：

第一章，绪论。本章首先介绍本文的研究背景和意义，阐述面向化学实验的具身智能技术面临的三个问题。随后从三维场景生成、三维物体重建以及面向化学实验的具身智能任务规划三个方面综述相关研究现状，并总结当前面临的主要挑战。最后给出本文的主要研究内容与章节安排。

第二章，理论基础。本章介绍本文后续研究所涉及的相关理论与关键技术基础，主要包括大语言模型基础、遗传算法、神经隐式场基础和本文实验中使用的相关评价指标，为后续章节的算法设计与实验分析提供理论支撑。

第三章，基于大语言模型的化学实验场景生成。本章针对化学实验场景在安全性、场景质量与多样性方面的需求，提出一种由大语言模型驱动的化学实验场景生成框架。首先，利用大语言模型从自然语言实验描述中生成物体清单与实验步骤，并结合化学知识库进行物体安全补齐，在此基础上构建物体约束图并生成初始场景。随后，设计面向场景展示的相机规划算法，以提高场景浏览与人工评审效率。最后，引入人机协同的场景进化算法，对生成场景进行迭代优化。通过场景生成实验、相机规划实验和进化算法实验，验证所提方法在安全性、质量与多样性方面的有效性。

第四章，基于隐式神经场的高保真化学资产重建。本章围绕化学实验资产类型多样且现有数据集覆盖不足的问题，首先提出了化学实验物体分类体系图，并据此分析不同类别物体带来的重建挑战。在此基础上，提出了一种基于隐式神经场 NeuS 的高保真化学资产重建算法，该算法能够有效应对各个类别物体所带来的重建挑战。随后，针对小物体提出学习态增强的重建加速策略，针对透明物体通过绝对距离场的局部极小值实现统一提取。在此基础上，进一步研究从隐式颜色场到高分辨率纹理贴图的自动化生成方法，实现从多视角图像到显式三维资产的完整转换。最后，通过透明物体几何重建、小物体重建与纹理贴图实验，验证所提方法在几何质量与纹理质量上的有效性。

第五章，面向化学实验的具身智能任务规划。本章围绕具身智能任务规划问题，提出一种“实验手册文本—步骤计划—操作原语”的分层转化方法。首先，利用大语言模型解析自然语言实验手册，生成结构化步骤计划；随后，结合预定义的机器人操作原语，将高层实验步骤进一步转化为机器人可调用的操作序列。最后，在基于 NVIDIA Omniverse 的机器人仿真平台中，以典型无机化学反应任务为例，展示从实验描述到机器人执行的完整流程，并通过生成质量与生成效率实验验证所提方法的有效性。

第六章，总结与展望。本章对全文工作进行总结，归纳本文在化学实验场景生成、高保真化学资产重建以及机器人任务规划方面的主要研究内容与结论，并结合当前方法的局限性，对后续可能的研究方向进行展望。